\documentclass{article}
\usepackage[utf8]{inputenc}
\usepackage{a4wide}

\begin{document}

\section*{สำรวจภูเขา}

กำหนดอาร์เรย์\texttt{arr}ของเลขจำนวนเต็มมาให้

ให้สำรวจหาภูเขาจากอาร์เรย์ย่อย\texttt{(subarray)}ที่มีเงื่อนไขดังต่อไปนี้



1. ความยาวอาร์เรย์ย่อยไม่น้อยกว่า 3



2. มีดัชนี\texttt{i}(ดัชนีเริ่มต้นจาก 0)

โดยที่\texttt{0\ \textless{}i\textless{}\ ความยาวอาร์เรย์\ –\ 1}และ{}\texttt{arr{[}0{]}\textless{}\ arr{[}1{]}\textless{}\ ...\ \textless{}\ arr{[}i–1{]}\textless{}\ arr{[}i{]}}{}และ{}\texttt{arr{[}i{]}\ \textgreater{}\ arr{[}i+1{]}\textgreater{}...\textgreater{}\ arr{[}arr.length\ –\ 1{]}}{}



อาร์เรย์ของจำนวนเต็ม\textbf{ให้หาความยาวของอาร์เรย์ย่อยภูเขาที่ยาวที่สุด

ตอบ\texttt{0}หากไม่มีอาร์เรย์ย่อยภูเขานั้น}\\



\textbf{\hfill\break

}



\textbf{ตัวอย่าง}



{\def\LTcaptype{none} % do not increment counter

\begin{longtable}[]{@{}

  >{\raggedright\arraybackslash}p{(\linewidth - 4\tabcolsep) * \real{0.3333}}

  >{\raggedright\arraybackslash}p{(\linewidth - 4\tabcolsep) * \real{0.3333}}

  >{\raggedright\arraybackslash}p{(\linewidth - 4\tabcolsep) * \real{0.3333}}@{}}

\toprule\noalign{}

\begin{minipage}[b]{\linewidth}\raggedright

ข้อมูลเข้า

\end{minipage} & \begin{minipage}[b]{\linewidth}\raggedright

ข้อมูลออก

\end{minipage} & \begin{minipage}[b]{\linewidth}\raggedright

คำอธิบาย

\end{minipage} \\

\midrule\noalign{}

\endhead

\bottomrule\noalign{}

\endlastfoot

\begin{minipage}[t]{\linewidth}\raggedright

7\\

2 1 4 7 3 2 5\strut

\end{minipage} & 5 & 1 4 7 3 2 \\

\begin{minipage}[t]{\linewidth}\raggedright

14\\

100 18 45 11 36 56 56 57 72 54 10 57 83 45\\

\strut

\end{minipage} & 5 & \begin{minipage}[t]{\linewidth}\raggedright

56 57 72 54 10 ฐาน 56 ซ้ำ ใช้ได้แค่ตัวเดียว\\

\strut

\end{minipage} \\

\begin{minipage}[t]{\linewidth}\raggedright

6\\

3 5 7 7 5 2\\

\strut

\end{minipage} & 0 & \begin{minipage}[t]{\linewidth}\raggedright

7 ซ้ำกันสองตัวไม่ได้เพราะภูเขาต้องยอดแหลม\\

\strut \\

\strut

\end{minipage} \\

\end{longtable}

}



\textbf{\hfill\break

}



\subsection*{Input}
บรรทัดแรก จำนวนเต็ม\texttt{N}เป็นขนาดของอาร์เรย์\texttt{arr}



บรรทัดที่สอง เป็นชุดเลขจำนวนเต็มแทนความสูง ณ แต่ละจุด แต่ละค่าจะเว้นด้วยช่องว่าง



\subsection*{Output}
จำนวนเต็มแทนความยาว\ul{\textbf{สูงสุด}}ของอาร์เรย์ย่อยภูเขา\\



\end{document}
